\chapter{Execution Engines}

\section{Two Engines, One Platform}

Connector provides two distinct execution engines, designed for complementary
problem spaces. Understanding when to use each one — and how they interoperate
— is central to building effective agents on the platform.

\begin{table}[H]
\centering
\renewcommand{\arraystretch}{1.5}
\begin{tabularx}{\textwidth}{@{}lXX@{}}
\toprule
 & \textbf{Cognitive Substrate} & \textbf{CLS Runtime} \\
\midrule
\textbf{What it is} &
  An 11-layer deliberative reasoning pipeline inspired by cognitive
  architectures (BDI, ACT-R, Soar) &
  A contract compiler and VM for governed, structured agent logic \\
\textbf{Best for} &
  Open-ended reasoning, planning under uncertainty, multi-step
  deliberation &
  Deterministic workflows, compliance enforcement, auditable
  structured execution \\
\textbf{Input} &
  Raw perception (text, sensor, multimodal) &
  A compiled SolutionContract (CID-addressed, signed) \\
\textbf{Output} &
  Commitment + action plan + reflection record &
  Signed ExecutionReceipt with full operation trace \\
\textbf{Governance} &
  Checkpointed at each of the 11 layers &
  Pre/post/invariant predicates enforced by the VM \\
\bottomrule
\end{tabularx}
\caption{Cognitive Substrate vs.\ CLS Runtime — when to use each engine}
\end{table}

Both engines write all intermediate and final state to the Memory Kernel as
MemPackets, ensuring that every reasoning step and every contract execution is
auditable and recoverable.

\section{Cognitive Substrate}

The Cognitive Substrate is Connector's reasoning engine. It transforms raw
inputs — text, sensor data, API responses — into committed, explainable actions
through an eleven-layer pipeline. Each layer produces a typed output that
becomes the input of the next, and each output is stored as a MemPacket with
its own CID.

\begin{insightbox}{Why a pipeline, not a loop?}
Most LLM agent frameworks use a simple observe → think → act loop. This works
for simple tasks but breaks down when agents need to explain their reasoning,
recover from mid-task failures, or satisfy compliance requirements. The 11-layer
pipeline makes each reasoning \emph{stage} explicit, checkpointed, and
independently verifiable.
\end{insightbox}

\subsection{The 11-Layer Thought Pipeline}

\begin{figure}[H]
\centering
\begin{tikzpicture}[
  step/.style={rectangle, rounded corners=3pt,
               draw=#1!70, fill=#1!10,
               text width=2.0cm, align=center,
               minimum height=1.05cm,
               font=\scriptsize\bfseries, inner sep=4pt},
  node distance=0.35cm and 0.35cm
]

%% Row 1 (layers 1-6)
\node[step=CKernel]  (l1) at (0,2.6)    {1\\Perception};
\node[step=CKernel]  (l2) at (2.45,2.6) {2\\Meaning};
\node[step=CExec]    (l3) at (4.9,2.6)  {3\\Tension\\Detection};
\node[step=CExec]    (l4) at (7.35,2.6) {4\\Knowledge\\Retrieval};
\node[step=CLogic]   (l5) at (9.8,2.6)  {5\\Possibility\\Generation};
\node[step=CLogic]   (l6) at (12.25,2.6){6\\Expertise\\Application};

%% Row 2 (layers 7-11, right to left so arrow flows left-to-right visually)
\node[step=CProto]   (l7)  at (12.25,1.1){7\\Evaluation};
\node[step=CProto]   (l8)  at (9.8,1.1)  {8\\Commitment};
\node[step=CDX]      (l9)  at (7.35,1.1) {9\\Planning};
\node[step=CDX]      (l10) at (4.9,1.1)  {10\\Reflection};
\node[step=CBiz]     (l11) at (2.45,1.1) {11\\Exposure};

%% Arrows row 1
\foreach \a/\b in {l1/l2, l2/l3, l3/l4, l4/l5, l5/l6}
  \draw[arr=gray!70] (\a) -- (\b);

%% Turn arrow
\draw[arr=gray!70] (l6.south) -- ++(0,-0.3) -- (l7.north |- 0,-0.0) -- (l7.north);

%% Arrows row 2 (right to left)
\foreach \a/\b in {l7/l8, l8/l9, l9/l10, l10/l11}
  \draw[arr=gray!70] (\a) -- (\b);

%% Labels
\node[above=0.15cm of l1, font=\scriptsize\itshape, text=gray] {raw input};
\node[below=0.15cm of l11, font=\scriptsize\itshape, text=gray] {action + explanation};

%% Checkpoint annotation
\draw[decorate, decoration={brace,amplitude=5pt,raise=3pt}, CKernel!50]
  (l1.north west) -- (l6.north east)
  node[midway, above=8pt, font=\scriptsize\itshape, text=CKernel]
    {MemPacket written at each layer};

\end{tikzpicture}
\caption{Cognitive Substrate 11-layer thought pipeline — each step is checkpointed as a MemPacket}
\label{fig:cognitive-pipeline}
\end{figure}

\begin{table}[H]
\centering
\renewcommand{\arraystretch}{1.35}
\begin{tabularx}{\textwidth}{@{}clXl@{}}
\toprule
\textbf{\#} & \textbf{Layer} & \textbf{What it does} & \textbf{Output type} \\
\midrule
1  & Perception           & Normalises raw input (text, sensor, multimodal) into a typed object & \texttt{PerceptionObject} \\
2  & Meaning Extraction   & Extracts entities, relations, intent from the perception & \texttt{MeaningObject} \\
3  & Tension Detection    & Identifies unresolved pressures — goals vs.\ reality gaps & \texttt{Tension[]} \\
4  & Knowledge Retrieval  & Queries memory for relevant facts (8 knowledge forms) & \texttt{KnowledgeSet} \\
5  & Possibility Generation & Generates candidate actions from tensions + knowledge & \texttt{Possibility[]} \\
6  & Expertise Application & Domain kernel scores each possibility (medical, finance…) & \texttt{EvaluationScores} \\
7  & Evaluation           & Multi-dimensional scoring: feasibility, impact, risk, cost & \texttt{RankedPossibilities} \\
8  & Commitment           & Selects best possibility; records revision history & \texttt{Commitment} \\
9  & Planning             & Builds a DAG of operations to execute the commitment & \texttt{ExecutionPlan} \\
10 & Reflection           & Compares expected vs.\ actual outcomes; extracts learnings & \texttt{ReflectionRecord} \\
11 & Exposure             & Renders result for human, peer agent, auditor, or debugger & \texttt{ExposurePacket} \\
\bottomrule
\end{tabularx}
\caption{Cognitive Substrate layer descriptions}
\end{table}

\subsection{Expertise Kernels}

Layer 6 is domain-aware. Rather than evaluating possibilities generically, the
Cognitive Substrate delegates to a pluggable \textbf{ExpertiseKernel} that
applies domain-specific rules. Built-in kernels include:

\begin{itemize}[leftmargin=1.5em, itemsep=2pt]
  \item \textbf{General} — Default scoring without domain constraints
  \item \textbf{Medical} — HIPAA-aware; flags PHI access, requires consent verification
  \item \textbf{Financial} — Compliance-aware; checks regulatory limits and disclosure rules
  \item \textbf{Legal} — Privilege-aware; enforces attorney-client boundary rules
\end{itemize}

Custom expertise kernels implement a single trait and are registered at boot time.

\section{CLS: Connector Logic Standard}

CLS is the \emph{contract system} of Connector: a language for declaring agent
logic, a compiler that transforms declarations into governed executables, and a
runtime VM that executes them with full predicate enforcement.

The key insight of CLS is that agent logic should be declared, compiled, and
verified \emph{before} it runs — not interpreted at runtime from prompts or
configuration files. A compiled contract is a content-addressed, signed
artifact: it cannot be tampered with, and any execution of it is fully
attributable.

\subsection{CLS Compiler Pipeline}

\begin{figure}[H]
\centering
\begin{tikzpicture}[
  pass/.style={rectangle, rounded corners=4pt,
               draw=CExec!70, fill=CExec!10,
               text width=2.8cm, align=center,
               minimum height=1.1cm,
               font=\footnotesize\bfseries, inner sep=5pt},
  io/.style={rectangle, rounded corners=4pt,
             draw=gray!50, fill=gray!8,
             text width=2.2cm, align=center,
             minimum height=0.9cm, font=\footnotesize},
  node distance=0.5cm and 0.55cm
]

\node[io] (src) {YAML surface\\declaration};
\node[pass, right=of src]  (p1) {1. Parse\\YAML → AST};
\node[pass, right=of p1]   (p2) {2. Desugar\\Normalise};
\node[pass, right=of p2]   (p3) {3. Lower\\AST → IR};
\node[pass, below=0.8cm of p3] (p4) {4. Verify\\Type check};
\node[pass, left=of p4]    (p5) {5. Optimise\\DCE, fold};
\node[pass, left=of p5]    (p6) {6. Emit\\IR → Contract};
\node[io, left=of p6] (out) {SolutionContract\\(CID-signed)};

\draw[arr=gray] (src) -- (p1);
\draw[arr=CExec!60] (p1) -- (p2);
\draw[arr=CExec!60] (p2) -- (p3);
\draw[arr=CExec!60] (p3) -- (p4);
\draw[arr=CExec!60] (p4) -- (p5);
\draw[arr=CExec!60] (p5) -- (p6);
\draw[arr=gray] (p6) -- (out);

\node[below=0.15cm of p4, font=\scriptsize\itshape, text=gray]
  {reachability, predicate sat.};
\node[below=0.15cm of p5, font=\scriptsize\itshape, text=gray]
  {dead state elim., const fold};

\end{tikzpicture}
\caption{CLS 6-pass compiler pipeline — YAML declaration to signed SolutionContract}
\label{fig:cls-compiler}
\end{figure}

Pass 4 (Verify) is the most important: the compiler checks that all predicates
are satisfiable, all state transitions are reachable, and all tool requirements
can be resolved. A contract that fails verification is \emph{never} stored and
\emph{never} executed.

\subsection{SolutionContract Anatomy}

\begin{figure}[H]
\centering
\begin{tikzpicture}[
  field/.style={rectangle, draw=CExec!50, fill=CExec!6,
                text width=4.0cm, align=left,
                minimum height=0.8cm, inner sep=5pt, font=\footnotesize},
  node distance=0.28cm
]

\node[rectangle, rounded corners=5pt, draw=CExec, fill=CExec!4,
      minimum width=10cm, minimum height=6.6cm,
      label={[font=\bfseries\small, text=CExec]above:SolutionContract}]
  (box) at (0,0) {};

\node[field] (cid)  at (0, 2.4) {\texttt{cid}\quad\footnotesize\itshape cls1-sha256-… (content address)};
\node[field] (iface) at (0, 1.6){\texttt{interface}\quad\footnotesize\itshape inputs, outputs, tool reqs, memory reqs};
\node[field] (sm)   at (0, 0.8) {\texttt{state\_machine}\quad\footnotesize\itshape states, transitions, operations};
\node[field] (gov)  at (0, 0.0) {\texttt{governance}\quad\footnotesize\itshape preconditions, postconditions, invariants};
\node[field] (res)  at (0,-0.8) {\texttt{resource\_envelope}\quad\footnotesize\itshape max tokens, memory, time, tools};
\node[field] (bill) at (0,-1.6) {\texttt{billing\_spec}\quad\footnotesize\itshape cost model, metering unit, rate};
\node[field] (sig)  at (0,-2.4) {\texttt{signature}\quad\footnotesize\itshape Ed25519 over canonical JSON};

\end{tikzpicture}
\caption{SolutionContract field anatomy — every field is included in the CID computation}
\label{fig:contract-anatomy}
\end{figure}

\subsection{Runtime Execution Lifecycle}

When a contract is executed, the CLS runtime follows a deterministic
seven-phase lifecycle:

\begin{figure}[H]
\centering
\begin{tikzpicture}[
  ph/.style={rectangle, rounded corners=3pt,
             draw=CExec!70, fill=CExec!9,
             text width=1.9cm, align=center,
             minimum height=0.95cm, font=\scriptsize\bfseries},
  node distance=0.4cm
]
\node[ph] (load)  {Load\\by CID};
\node[ph, right=of load]  (val)  {Validate\\signature};
\node[ph, right=of val]   (res)  {Resolve\\tools \& mem};
\node[ph, right=of res]   (init) {Initialise\\state machine};
\node[ph, right=of init]  (exec) {Execute\\with predicates};
\node[ph, right=of exec]  (rec)  {Emit\\Receipt};
\node[ph, right=of rec]   (clean){Cleanup\\resources};

\foreach \a/\b in {load/val, val/res, res/init, init/exec, exec/rec, rec/clean}
  \draw[arr=CExec!60] (\a) -- (\b);

\node[below=0.25cm of exec, font=\scriptsize\itshape, text=gray]
  {11-layer predicate chain};
\node[below=0.25cm of rec, font=\scriptsize\itshape, text=gray]
  {signed, CID-anchored};

\end{tikzpicture}
\caption{CLS runtime execution lifecycle — deterministic seven-phase flow}
\label{fig:cls-runtime}
\end{figure}

\section{Resource Budgets}

Both engines enforce a shared resource budget model. Budgets are specified when
a session is created and enforced by the kernel in real time. Budget violations
trigger immediate session termination, and the partial execution is still
receipted — so operators can see exactly how much was consumed before the limit
was hit.

\begin{table}[H]
\centering
\renewcommand{\arraystretch}{1.4}
\begin{tabularx}{\textwidth}{@{}llX@{}}
\toprule
\textbf{Dimension} & \textbf{Unit} & \textbf{What it controls} \\
\midrule
\texttt{max\_tokens}       & tokens         & Total LLM tokens consumed in the session \\
\texttt{max\_memory\_bytes}& bytes          & Total MemPacket storage written \\
\texttt{max\_wall\_time\_ms}& milliseconds  & End-to-end session wall-clock time \\
\texttt{max\_tool\_calls}  & count          & Number of external tool invocations \\
\texttt{max\_memory\_ops}  & count          & Number of kernel read/write syscalls \\
\bottomrule
\end{tabularx}
\caption{Resource budget dimensions enforced by both execution engines}
\end{table}

\begin{warnbox}{Budget violation behaviour}
When any budget dimension is exhausted, the session is terminated immediately.
The kernel writes a final \texttt{BudgetExceeded} receipt with the exact
consumption at time of violation. No partial results are silently discarded —
the partial execution is fully traceable in the receipt chain.
\end{warnbox}

\section{Engine Integration Patterns}

The two engines are designed to complement each other. Complex agent deployments
typically use both:

\begin{figure}[H]
\centering
\begin{tikzpicture}[
  eng/.style={rectangle, rounded corners=5pt, draw=#1!70, fill=#1!8,
              text width=4.2cm, align=center, minimum height=1.4cm,
              font=\small\bfseries},
  node distance=2.2cm
]
\node[eng=CKernel] (cog) {Cognitive Substrate\\{\footnotesize deliberation, planning}};
\node[eng=CExec, right=3.5cm of cog] (cls) {CLS Runtime\\{\footnotesize structured execution}};

\draw[arr=CKernel!60, bend left=20]
  (cog.north) to node[above, font=\footnotesize\itshape]
    {invoke contract for structured step} (cls.north);
\draw[arr=CExec!60, bend left=20]
  (cls.south) to node[below, font=\footnotesize\itshape]
    {delegate open-ended reasoning} (cog.south);

\node[below=2.0cm of cog, font=\footnotesize, text=gray, align=center,
      text width=4cm]
  {\textit{Use when:}\\open-ended, uncertain,\\multi-step reasoning needed};
\node[below=2.0cm of cls, font=\footnotesize, text=gray, align=center,
      text width=4cm]
  {\textit{Use when:}\\deterministic, auditable,\\compliance-critical workflow};

\end{tikzpicture}
\caption{Engine integration — Cognitive and CLS engines delegate to each other as needed}
\label{fig:engine-integration}
\end{figure}
